\section{Governing theorem of the radical lattice and the power--vacancy law}

The appearance of \(369\) and \(693\) is not reducible to a graphic similarity. There are two exact mechanisms that converge:

\begin{enumerate}
\item \textbf{nonadic algebraic mechanism:} \(3,6,9\) are the visual representatives of the nilpotent radical \(\mathfrak m=(3)\subset\mathbb Z/9\mathbb Z\); the words \(369\), \(693\), and \(936\) are the three phases of its multiplicative orbit;
\item \textbf{Archimedean vacancy mechanism:} the difference between the decimal length of \(9^n\) and the boundary \(10^n\) is governed exactly by the same irrational rotation \(\delta_{\rm vac}=\log_{10}(10/9)\) that already organizes the vacancy calendar.
\end{enumerate}

For \(n=9^9\), both mechanisms meet in the identity

\[
9^9-\left\lceil 9^9\log_{10}\frac{10}{9}\right\rceil
=369\,693\,099,
\]

so that

\[
\#\operatorname{cifras}\bigl(9^{9^9}\bigr)=369\,693\,100.
\]

The block \(369\mid693\) is, therefore, an exact output of the logarithmic reader of the nonadic power; it is not the decimal prefix of \(9^{9^9}\). The final unit proceeds from the universal operator that transforms decimal order into decimal length.

This identity opens a broader structural reading: addition, product, power, subtraction, division, root, and logarithm can be organized as a hierarchy of transports and quotients. Each ascent of operation compactifies histories; TPK preserves the coordinates that make their reconstruction possible.


\section{Radical structure of the APP on \(\mathbb Z/9\mathbb Z\)}

\subsection{Nonadic local ring \(\mathbb Z/9\mathbb Z\) and nilpotent radical}

Let

\[
R=\mathbb Z/9\mathbb Z.
\]

Its group of units and its maximal ideal are

\[
R^\times=\{1,2,4,5,7,8\},
\qquad
\mathfrak m=(3)=\{0,3,6\}=\{9,3,6\}_{\rm vis}.
\]

Moreover,

\[
\mathfrak m^2=0.
\]

The preceding equality is an equality of ideals: every product of two elements of \(\mathfrak m\) vanishes modulo nine. It must not be confused with the Cartesian product \(\mathfrak m\times\mathfrak m\), which contains nine APP positions.

\subsection{Exact sequence of the multiplier \(3\)}

Multiplication by three,

\[
M_3:R\longrightarrow R,
\qquad x\longmapsto3x,
\]

has

\[
\ker M_3=\mathfrak m,
\qquad
\operatorname{im}M_3=\mathfrak m.
\]

It therefore induces an isomorphism of spaces over \(\mathbb F_3\),

\[
\overline M_3:R/\mathfrak m\xrightarrow{\sim}\mathfrak m,
\qquad
\overline x\longmapsto3x.
\]

In visual representatives,

\[
\overline0\longmapsto9,
\qquad
\overline1\longmapsto3,
\qquad
\overline2\longmapsto6.
\]

Thus, the band \(3,6,9\) is not a decimal ornament. It is an internal copy of the ternary quotient lodged in the nonadic radical.

\subsection{Cyclic orbit \(369\to693\to936\to369\)}

The row of the multiplication table corresponding to \(3\) is

\[
3,6,9,3,6,9,3,6,9.
\]

The cambio of phase \(j\mapsto j+1\) produces

\[
369\longmapsto693\longmapsto936\longmapsto369.
\]

These three words are the three cyclic readings of the application.

\[
\mathbb F_3\longrightarrow\mathfrak m,
\qquad
\bar j\longmapsto3j.
\]

The row \(6\), by satisfying \(6\equiv-3\pmod9\), reverses the orientation
of the same cycle. The mark \(9\) fulfills two typed functions:

\[
9\bmod9=0,
\qquad
\operatorname{dr}_9(9)=9.
\]

The module publishes the null class; the positive digital root preserves the mark of closure. This difference of readers is precisely what allows \(9\) to be at once residual zero and visible closure.

\subsection{Radical cross of APP and coordinate reduction to \(\mathbb F_3^2\)}

On

\[
X_9=R^2,
\]

the union

\[
\mathscr C_{369}
=(\mathfrak m\times R)\cup(R\times\mathfrak m)
\]

is the cross formed by rows and columns \(3,6,9\). It contains

\[
3\cdot9+3\cdot9-3\cdot3=45
\]

positions. Their central intersection \(\mathfrak m\times\mathfrak m\) contains nine positions. In the multiplicative observable, the sum of the cross is

\[
297=216+81=3(72+27).
\]

The coordinate reduction produces

\[
X_9/(\mathfrak m\times\mathfrak m)\cong\mathbb F_3^2
\]

as a quotient set typed by the classes modulo three. The inverse lift requires preserving the coordinate within each three-element fiber.

\subsection{Kernel \(7227\), spectral decomposition and link to the radical cross}

The central block of the APP is

\[
B_c=
\begin{pmatrix}
7&2\\
2&7
\end{pmatrix}
=7I+2R
=9P_++5P_-.
\]

Its rowwise reading produces \(72\mid27\); its diagonal reading produces \(77\mid22\). One verifies

\[
72+27=77+22=99,
\]

\[
72-27=45=\det B_c,
\]

\[
77-22=55=54+1.
\]

The relation with the radical cross is not nominal:

\[
3\cdot72=216,
\qquad
3\cdot27=81,
\qquad
3\cdot99=297.
\]

The central unitary block and the non-unitary network \(3,6,9\) are thereby linked by an exact local--global identity. The hierarchy of the sum--product difference preserves its types:

\[
4\quad\text{(local spectral jump)},\qquad
2\quad\text{(off-diagonal coupling)},
\]

\[
6\quad\text{(compressed difference modulo three)},\qquad
54\quad\text{(two-dimensional unfolding)}.
\]


\section{Power--vacancy law for the scales \texorpdfstring{\(9^n\)}{9 to the power n} and \texorpdfstring{\(10^n\)}{10 to the power n}}

\subsection{Decimal Length Functions and Accumulated Vacancy}

For \(n\ge1\), let

\[
D_9(n)=\#\operatorname{cifras}(9^n)
=\left\lfloor n\log_{10}9\right\rfloor+1.
\]

The decimal boundary \(10^n\) has \(n+1\) digits. We define the accumulated vacancy of the nonadic power by

\[
V_9(n)=(n+1)-D_9(n).
\]

We introduce the two complementary numbers

\[
\rho_{\rm vac}=\log_{10}9,
\qquad
\delta_{\rm vac}=1-\rho_{\rm vac}
=\log_{10}\frac{10}{9}.
\]

\subsection{Exact identity between decimal length and vacancy rotation}

\textbf{Theorem.} For every integer \(n\ge1\),

\[
\boxed{V_9(n)=\left\lceil n\delta_{\rm vac}\right\rceil},
\]

and, therefore,

\[
\boxed{D_9(n)=n-\left\lceil n\delta_{\rm vac}\right\rceil+1}.
\]

\textbf{Proof.} Since \(\delta_{\rm vac}\) is irrational, \(n\delta_{\rm vac}\notin\mathbb Z\). Then

\[
\begin{aligned}
D_9(n)
&=\lfloor n(1-\delta_{\rm vac})\rfloor+1\\
&=\lfloor n-n\delta_{\rm vac}\rfloor+1\\
&=n-\lceil n\delta_{\rm vac}\rceil+1.
\end{aligned}
\]

Subtracting from \(n+1\) yields the first identity. \(\square\)

\subsection{Mechanical coding of vacancies through nonadic powers}

The increments of the vacancy are

\[
Z_n=V_9(n+1)-V_9(n)
=\lceil(n+1)\delta_{\rm vac}\rceil
-\lceil n\delta_{\rm vac}\rceil
\in\{0,1\}.
\]

This is exactly the mechanical word of slope \(\delta_{\rm vac}\) used by the HMT vacancy calendar. Therefore, that calendar admits an additional, completely exact reading:

\begin{quote}
\(Z_n=1\) if and only if, in passing from \(9^n\) to \(9^{n+1}\), a new decimal range is lost relative to the boundary \(10^n\to10^{n+1}\).
\end{quote}

The power does not add a coincidence to the theory of vacancies: it provides its natural realization as a counter of decimal rank loss.

\subsection{Law \(21/22\) as a case of the power--vacancy theorem}

Since

\[
\delta_{\rm vac}^{-1}
=21.85434532678283256\ldots,
\]

the indices for which \(Z_n=1\) are separated by \(21\) or \(22\). The first twenty-one exponents satisfy

\[
D_9(n)=n,\qquad1\le n\le21,
\]

whereas

\[
D_9(22)=21.
\]

In particular,

\[
9^9=387\,420\,489
\]

has exactly nine digits. This initial preservation of rank is not an indefinite rule: the vacancy calendar determines exactly when it ceases to hold.

\subsection{Exact evaluation at \(n=9^9\) and length of \(9^{9^9}\)}

Let

\[
n_*=9^9=387\,420\,489.
\]

Then

\[
n_*\delta_{\rm vac}
=17\,727\,389.3684296412564569\ldots,
\]

and, therefore,

\[
\boxed{V_9(n_*)=17\,727\,390}.
\]

The conserved coordinate after discounting vacancies is

\[
n_*-V_9(n_*)
=387\,420\,489-17\,727\,390
=369\,693\,099.
\]

Therefore,

\[
\boxed{\operatorname{ord}_{10}\bigl(9^{9^9}\bigr)=369\,693\,099},
\]

\[
\boxed{D_9(9^9)=369\,693\,100}.
\]

Furthermore, the congruences are obtained.

\[
9^9\equiv0\pmod9,\qquad
V_9(9^9)\equiv0\pmod9,
\]

\[
369\,693\,099\equiv0\pmod9,\qquad
369\,693\,100\equiv1\pmod9.
\]

The corresponding digit sums are

\[
s_{10}(9^9)=45,\qquad
s_{10}(V_9(9^9))=36,
\]

\[
s_{10}(369\,693\,099)=54,\qquad
s_{10}(369\,693\,100)=37.
\]

The return \(9\to1\) appears here with precise types:

\begin{enumerate}
\item \(9^9\) and the accumulated vacancy belong to the null class modulo nine;
\item its difference, the decimal order, remains in the null class and has digit sum \(54\);
\item the universal operator \(D=\operatorname{ord}+1\) adds the closure unit and publishes the positive class \(1\).
\end{enumerate}

\subsection{Domain of validity of the power--vacancy law}

It is proved:

\begin{itemize}
\item the theorem potencia–vacancia for every \(n\);
\item the identity with the mechanical vacancy word;
\item the exact evaluation in \(n=9^9\);
\item the decimal order and length of \(9^{9^9}\);
\item the preceding congruences and digit sums.
\end{itemize}

It is not claimed that \(9\) is the only basis with an analogous property in every domain, that each external occurrence of \(369\) arises from this tower, or that the tower replaces the construction of the nine phases or the generation certificate for \(\pi,e,\varphi\).


\section{Operator Hierarchy, Reversibility, and Genealogical Memory}

\subsection{Translation, Dilation, and Exponentiation as Operational Levels}

Elementary operations can be seen as actions:

\[
T_a(x)=x+a,\qquad
M_a(x)=ax,\qquad
P_k(x)=x^k.
\]

Their partial inverses are, respectively, subtraction, division, and extraction of roots. The passage from one operation to the next modifies the structure of the fibers.

\subsection{Stratified reversibility in \(\mathbb Z/9\mathbb Z\)}

\textbf{proposition.} En the anillo nonadico:

\begin{enumerate}
\item every translation \(T_a\) is bijective;
\item \(M_a\) is bijective if and only if \(a\in R^\times\);
\item if \(a\in\mathfrak m\), multiplication has nontrivial fibers and
its image lies within \(\mathfrak m\);
\item for \(x\in\mathfrak m\), \(x^k=0\) for every \(k\ge2\);
\item over \(R^\times\cong C_6\), the map \(P_k\) is an automorphism if and only if \(\gcd(k,6)=1\).
\end{enumerate}

\textbf{Proof.} The first assertion uses the additive group structure. The second is the characterization of the units. The third follows from \(aR\subseteq\mathfrak m\). The fourth follows from \(\mathfrak m^2=0\). The fifth is the condition for multiplication by \(k\) to be an automorphism of the cyclic group of order six. \(\square\)

The result produces an exact hierarchy:

\[
\text{sum: global reversibility},
\]

\[
\text{product: reversibility restricted to units},
\]

\[
\text{power: cyclic dynamics on units and nilpotent collapse in }\mathfrak m.
\]

Roots are functions only when the corresponding power is injective on the chosen domain; on the total domain they are multibranch correspondences or empty relations.

\subsection{Reading of the Topological Phase Kernel on the operator hierarchy}

This stratification suggests a functional definition:

\begin{quote}
TPK does not add an external operation to sum, product, and power. It preserves the sheet, quotient, carry, orientation, route, and boundary that each ordinary evaluation eliminates when identifying several histories with the same output.
\end{quote}

Sum and product are the first two evaluations on the APP. A power is repeated composition of the product; a root selects or reconstructs branches of that composition. The record distinguishes the published value from the operatorial history producing it.

\subsection{Logarithmic descent between consecutive levels of the hierarchy}

In a positive domain,

\[
\log(xy)=\log x+\log y,\qquad
\log(x^k)=k\log x,\qquad
\log(a^x)=x\log a.
\]

The logarithm converts product into sum, power into product, and exponential into dilation. The defect appears when the subsequent reader discretizes:

\[
\mathcal L_{10}(a^n)
=\lfloor n\log_{10}a\rfloor+1.
\]

The integer part introduces the boundary-crossing calendar. For \(a=9\), that calendar is precisely the vacancy word.

\subsection{Affine conjugation of translations and dilations}

Translations and dilations satisfy

\[
M_aT_bM_a^{-1}=T_{ab}
\]

when \(a\) is invertible. This identity expresses that the product reorganizes the units of addition. When \(a\) ceases to be invertible, conjugation breaks because branches disappear. The radical \(3,6,9\) localizes precisely that failure.


\section{Path Algebra: Multiplicative Composition, Additive Action and Superposition}

In a path algebra, addition aggregates alternatives and multiplication concatenates trajectories. For local weights \(w(e)\),

\[
K(x,y)
=\sum_{\gamma:x\to y}
\prod_{e\in\gamma}w(e).
\]

If a coherent logarithmic branch is fixed,

\[
\prod_{e\in\gamma}w(e)
=\exp\!\left(\sum_{e\in\gamma}\log w(e)\right),
\]

and, with the appropriate discrete action,

\[
K(x,y)=\sum_{\gamma:x\to y}e^{-S(\gamma)/\hbar}.
\]

The equation combines the operations:

\begin{itemize}
\item the product composes successive steps;
\item the logarithm accumulates the local action;
\item the exponential reconstructs the route weight;
\item the suma agrega all the alternativas.
\end{itemize}

APP is the minimal support on which sum and product act on the same histories. TPK adds the data that prevent the sum over paths from becoming a sum over values without genealogy.

In this framework, a digit is the output of a reader; a value is the Archimedean projection of a section; an HMT number preserves the compatible family of paths that supports that value; a particle preserves that genealogy under a persistent spectral realization.


\section{Internal root of \(-1\), exponential structure, and elliptic, parabolic, and hyperbolic geometries}

\subsection{Internal root of \(-1\) in the finite branch}

The Paley incidence matrix provides an operator \(A_W\) with

\[
A_W^2=-I.
\]

This turns \(\mathbb F_3^6\) into a three-dimensional space over

\[
\mathbb F_9=\mathbb F_3[\iota]/(\iota^2+1).
\]

The root of \(-1\) is not adjoined from outside as complex notation. It emerges as an endomorphism of an already constructed ternary support.

\subsection{Continuous realization of the induced quadratic type}

In the elliptic real branch, the generator \(J_{\rm e}\) satisfies

\[
J_{\rm e}^2=-I,
\]

and

\[
e^{tJ_{\rm e}}=\cos t\,I+\sin t\,J_{\rm e}.
\]

In particular,

\[
e^{\frac\pi2J_{\rm e}}=J_{\rm e},\qquad
e^{\pi J_{\rm e}}=-I,\qquad
e^{2\pi J_{\rm e}}=I.
\]

Thus, \(J_{\rm e}\) is the internal quarter-turn and Euler's equation is its half-turn closure.

\subsection{Common operator support of \(3\) geometric readings}

Let

\[
\mathbb D=\{z\in\mathbb C:|z|<1\}.
\]

This same set contains the real diameter \((-1,1)\), the complex
interior, the trigonometric boundary \(S^1\), and the Poincaré
hyperbolic metric

\[
ds_{\rm P}^2=\frac{4\ell^2|dz|^2}{(1-|z|^2)^2}.
\]

The map

\[
R_{\pi/2}(z)=iz
\]

preserves \(|z|\) and \(|dz|\). By consiguiente, is simultaneamente:

\begin{itemize}
\item complex multiplication by \(\sqrt{-1}\);
\item trigonometric rotation of \(90^\circ\);
\item Euclidean isometry of the disk;
\item Poincaré hyperbolic isometry;
\item automorphism of its circular boundary.
\end{itemize}

The coincidence is literal: the same map acts on the same carrier; the metric reader changes.

\subsection{Elliptic, parabolic, and hyperbolic geometries of the extended Euler equation}

The generators

\[
J_{+1}^2=-I,\qquad
J_0^2=0,\qquad
J_{-1}^2=I
\]

produce

\[
e^{tJ_{+1}}=\cos t\,I+\sin t\,J_{+1},
\]

\[
e^{tJ_0}=I+tJ_0,
\]

\[
e^{tJ_{-1}}=\cosh t\,I+\sinh t\,J_{-1}.
\]

The same exponential series separates into three regimes according to the
square of the generator: elliptic, parabolic, and hyperbolic. The
specialization

\[
e^{2\log\varphi\,J_{-1}}=F_{\rm av}^{\,2}
\]

integrates the Fibonacci self-scale. Euler and Fibonacci appear as
specializations of a common quadratic exponential family.

\subsection{Common carrier, common operator, and quadratic exponential family}

Geometric fusion is formulated at three levels:

\begin{enumerate}
\item \textbf{common carrier:} the unit disk contains the real diameter, the complex interior, and the circular boundary;
\item \textbf{common operator:} \(z\mapsto iz\) is a quarter-turn, complex multiplication, and a hyperbolic isometry;
\item \textbf{common family:} the tritic exponential equation brings together the elliptic, parabolic, and hyperbolic branches.
\end{enumerate}

The result does not identify the Euclidean and hyperbolic metrics. It states
that they are distinct realizations on the same support and share the
automorphism originating in the internal root of \(-I\).


\section{Restriction of internal rotation, edge holonomy and minimal torsion}

The construction contains two positive and differentiated objects:

\begin{enumerate}
\item the internal generator \(J_{\rm e}\), with \(J_{\rm e}^2=-I\);
\item the global invariant
\end{enumerate}
   \[
   \Delta_4=
   \frac{\pi-e\log\pi}{270}
   \left(1+\frac{\pi}{729}\right),
   \]
   denoted minimal torsion and used prior to species selection.

The authorial thesis adds that internal rotational dynamics, when published
at the boundary, is absorbed as minimal torsion. The relation must not be
written as a bare equality \(J_{\rm e}=\Delta_4\): the objects have different types.
Its natural form requires

\[
\text{interior rotation}
\xrightarrow{\operatorname{Res}_{\partial}}
\text{boundary holonomy}
\xrightarrow{\operatorname{Def}}
\text{scalar defect}.
\]

Let \(q_\partial\) be the boundary reader, \(\nabla_\partial\) the induced connection, and \(e_\partial\) the discrete frame. The edge torsion has type

\[
T_\partial=d_{\nabla_\partial}e_\partial.
\]

The author's claim acquires probative form through a functional

\[
\Lambda_\partial:
\operatorname{Tors}(\partial X)
\longrightarrow\mathbb R
\]

such that

\[
\Lambda_\partial(T_\partial)=\Delta_4
\]

and a theorem deriving \(T_\partial\) from the restriction of \(J_{\rm e}\) through the boundary operators of the TPK.

The internal root of \(-I\), the quarter-turn, the extended Euler equation, and the formula for \(\Delta_4\) are proved in their respective domains. The passage from interior rotation to minimal boundary torsion requires the declared morphism \(\operatorname{Res}_{\partial}\) and functional \(\Lambda_{\partial}\); it is not identified by a scalar equality.


\section{Generational triality, conjugation, and mixing over the enriched state}

The phrase of Isidor Isaac Rabi may serve as a historical reference, but it must not govern the title or constitute an autonomous section. The scientific question is:

\begin{quote}
What internal mechanism distinguishes three stable orders of fermionic realization and preserves that distinction under color, conjugation, and mixing?
\end{quote}

The HMT response is not "because TRIT has three values". The structure separates

\[
\operatorname{flav}(X)
=\bigl(\sigma_{ud}(X),g(X),\chi_{\rm fina}(X)\bigr),
\]

where \(\sigma_{ud}\) distinguishes branches, \(g\) distinguishes physical generations, and \(\chi_{\rm fina}\) preserves orientation and memory within a branch. For quarks, the color orbit is added after flavor and before the character.

Triplicity is obtained as a consequence of the enriched structure; its
effects on color, conjugation, and mixing are developed below:

\begin{enumerate}
\item color, flavor, generation, and mass are distinct readers of the same enriched state;
\item particle–antiparticle conjugation acts on the orientation without erasing the generational class;
\item CKM and PMNS transport already typed state bases; they do not create triplicity;
\item the neutrino classification must be derived from representation, occupation, and mixing, not by renaming depths \(G0,\ldots,G4\).
\end{enumerate}

That neutrinos are half-integer-spin leptons obeying fermionic statistics is a known physical classification. The structural contribution of HMT consists in obtaining internally why the neutral sector occupies that representation, why it admits three physical realizations, and how it couples to the mixing operator without reading its external angles.


\section{Operator memory of the TPK: preservation of residue, quotient, carry, sheet, orientation, route, and boundary under projection}

\subsection{Loss of reversibility and genealogical lifting}

In APP, different histories may produce the same residue. The product folds more histories than the sum; the power folds more than the product; the root attempts to reconstruct branches already identified. TPK preserves the datum that makes it possible to invert that loss:

\[
\text{residue},\quad
\text{quotient},\quad
\text{carry},\quad
\text{sheet},\quad
\text{orientation},\quad
\text{route},\quad
\text{boundary},\quad
\text{memory}.
\]

An algebraic operation may be characterized not only by its value, but also by the type of genealogy that it preserves or destroys.

\subsection{Projection of Compatible Histories onto the Archimedean Continuum}

The real line appears when a family of compatible cylinders contracts under Archimedean evaluation. The real value is the shadow of the section, not its complete genealogy. The exceptional projection preserves another part as incidence and symmetry.

The path integral expresses the same principle:

\[
\text{alternatives}\xrightarrow{\sum}\text{global amplitude},
\qquad
\text{successive stages}\xrightarrow{\prod}\text{route weight}.
\]

The logarithm converts multiplicative composition into additive action; the exponential returns it to amplitude. The discrete structure of the continuum can be formulated as genealogical conservation of all routes before their projection onto \(\mathbb R\).

\subsection{Distinction between numerical section and persistent realization of particle}

A number and a particle share minimal ontological invariants: memory, orientation, persistence, closure, signature, and position in a family of prolongations.

The number publishes an Archimedean coordinate. The particle adds spectral realization, occupation, propagation, and stability. The passage HMT \(\to\) MD changes the reader applied to the same enriched genealogy; it does not arbitrarily change the object.


\section{Joint closure of the radical network, the power--vacancy law, Euler geometry, and the operative hierarchy}

\subsection{Theorem of the Nonadic Radical Network}

The APP net \(3,6,9\) is the visual representation of the nilpotent radical \((3)\subset\mathbb Z/9\mathbb Z\). Its phases \(369\), \(693\), and \(936\) realize the isomorphism \(R/(3)\cong(3)\) induced by multiplication by three. The bands locate exactly the failure of invertibility of the product.

\subsection{Power-void Theorem}

For every \(n\ge1\), the number of decimal ranks lost by \(9^n\) relative to \(10^n\) is \(\lceil n\log_{10}(10/9)\rceil\). Their increments form the mechanical vacancy word. In \(n=9^9\), the retained coordinate is \(369\,693\,099\) and the resulting length is \(369\,693\,100\).

\subsection{Exponential Family Euler's Geometric Theorem}

The internal root of \(-I\) produces a quarter-turn. On the unit disk, the same map \(z\mapsto iz\) is simultaneously complex multiplication, trigonometric rotation, and a Poincaré isometry. The tritic exponential equation extends this regime to the parabolic and hyperbolic regimes, incorporating the Fibonacci self-scaling in the latter.

\subsection{Operatorial conservation principle between levels}

The sum--product--power hierarchy is a hierarchy of fibers and loss of
reversibility. The TPK is the lift that preserves the coordinates required
to distinguish histories identified by the ordinary readers.


\section{Typing hypotheses and negative controls}

\begin{enumerate}
\item Do not identify \(369\,693\,100\) with a TPK word without declaring the reader between domains.
\item No confundir the ideal \(\mathfrak m^2=0\) With \(\mathfrak m\times\mathfrak m\).
\item Do not confuse modulo nine with digital root: \(9\mapsto0\) in the former and \(9\mapsto9\) in the latter.
\item Do not use the threefold character of TRIT as a sufficient proof of three physical generations.
\item Do not turn \(J^2=-I\) into \(\Delta_4\) without the boundary map and the torsion functional.
\item Do not present \(369\) as a prefix of \(9^{9^9}\): it pertains to its order and its decimal length.
\item Do not reduce the ten thousand digits to an order-of-magnitude test: they certify extensive truncations of a parametrizable process whose thesis is indefinite prolongation.
\end{enumerate}

\section{Joint consequences for the radical net, the operational memory, and the edge torsion}

The principal finding is not that a tower visually contains the digits \(3,6,9\). The nonadic power and the vacancy calendar are two presentations of the same logarithmic defect, while APP localizes \(3,6,9\) as the radical where multiplication loses invertibility. The tower connects both extremes:

\[
\text{nonadic radical}
\longrightarrow
\text{power}
\longrightarrow
\text{logarithm}
\longrightarrow
\text{vacancy}
\longrightarrow
\text{decimal order}
\longrightarrow
\text{closure }9\to1.
\]

The extended Euler equation adds the second direction: the internal root of \(-I\) converts the discrete structure into rotational dynamics and gathers, on a common support, the real line, the trigonometric circle, and the hyperbolic disk. The later passage to minimal torsion must preserve the boundary map already developed in the corpus, not replace it with a nominal identification.
